\chapter{Limitations and Future Work}
\label{ch:future-work}

\section{Known Limitations}
\begin{itemize}
  \item \textbf{GPS accuracy.} Consumer GPS is accurate to roughly
        3--15\,m and degrades near buildings. Against a 10\,m geofence this
        can require a user standing at a checkpoint to wait for the fix to
        settle before the arrival prompt fires; the per-place radius
        mitigates but does not eliminate this.
  \item \textbf{Secure-context requirement.} Camera and precise
        geolocation APIs are available only over HTTPS (or localhost); on
        an insecure origin the system degrades to the compass-only HUD.
  \item \textbf{External media dependencies.} Waypoint media is served
        from Cloudinary and avatar images from a public CDN; both require
        connectivity, and there is no offline mode.
  \item \textbf{Compass variability.} Device magnetometers differ in
        quality; non-absolute headings trigger a calibration banner and
        may need a figure-eight gesture.
  \item \textbf{Permanent quiz attempts.} An attempt is permanent even if
        an administrator later edits the quiz content.
  \item \textbf{Single language.} The interface and content are
        English-only.
\end{itemize}

\section{Future Work}
Planned directions, in rough priority order: additional heritage routes
(the Kathmandu Valley UNESCO properties are natural candidates);
multi-language content, which matters more for heritage visitors than for
campus use; an in-dashboard quiz authoring interface (quizzes are currently
managed through the API and seed scripts); richer analytics such as dwell
time and per-waypoint drop-off; offline route bundles for areas with poor
connectivity; native app packaging for tighter sensor access; and an
accessibility pass across the exploration flow.
